% GENERATED FILE. Edit canonical lessons, then run npm run generate:books.
% canonical-source-hash: fnv1a64:c10f05f6c56854af
% canonical-lessons: ZH-C68-shan, ZH-W68-shan, ZH-C68-xingxing, ZH-C68-cao, ZH-C68-hua, ZH-C68-taiyang

\chapter{Out of Town}
\label{ch:zh-out-of-town}
\hlchaptermodality{68}

\begin{chapteropening}
\textbf{By the end of this chapter:} \emph{I can name mountains, stars, grass, flowers and the sun.}

Complete the last lesson of chapter 68: I can name mountains, stars, grass, flowers and the sun.
\end{chapteropening}

\section[\texorpdfstring{\zh{山} (shān)}{山 (shān)}]{\zh{山} (shān) — mountain; hill}

\label{lesson:ZH-C68-shan}



\begin{quote}
Before the new one: say the Mandarin for bird, then the Mandarin for fish.
\end{quote}

\subsection*{You'll want to know: \zh{山}}
\textbf{\zh{山}} — \emph{shān} — \textquotedblleft{}mountain; hill\textquotedblright{}.

\begin{quote}
\textbf{\zh{我喜欢山}} \emph{wǒ xǐhuan shān} — I love the mountains.
\end{quote}

The new character, in pieces: \zh{山} — three peaks standing on one base. The next lesson puts it in your hand.

\begin{cousinweb}
\textbf{\zh{山}} (\emph{shān}) is inherited from Middle Chinese \emph{srean} (Guangyun, in Baxter's transcription); the dictionary quotes an old book saying the mountain brings forth all things.
\end{cousinweb}

\begin{grammarlens}[title={What you've built}]
One more word for everyday answers.
\end{grammarlens}

\subsection*{Your turn}
\begin{itemize}
\raggedright
  \item \emph{Say it:} \emph{shān}
  \item \emph{Say it:} \emph{shān}, once more
  \item \emph{Say it:} \emph{wǒ xǐhuan shān} — I love the mountains.
  \item \emph{Recall:} say the Mandarin for bird, then the Mandarin for fish, then say \emph{shān} again
\end{itemize}

\begin{tcolorbox}[breakable,colback=teal!4,colframe=teal!35!black,title={Before you move on}]
What does \zh{山} mean? (\textquotedblleft{}Mountain; hill\textquotedblright{}.) Say it once more.
\end{tcolorbox}
\section[\texorpdfstring{\zh{山} (shān)}{山 (shān)}]{\zh{山} — three peaks on a base}

\label{lesson:ZH-W68-shan}



\begin{quote}
The word first, then the character inside it.

\begin{itemize}
\raggedright
  \item \emph{Recall:} say \emph{shān} — \textbf{R1}, one lesson back, before it decays
\end{itemize}

\begin{quote}
\zh{岁} \zh{中}
\end{quote}

\emph{Read it:} each one aloud, with nothing said first — \emph{suì}, \emph{zhōng}
\end{quote}

\subsection*{Script — three peaks on a base}
\hlblockfigure{\detokenize{figures/ZH-W68-shan-filmstrip.pdf}}{How \zh{山} is written, stroke by stroke}

\begin{quote}
\zh{山}
\end{quote}

\textbf{One piece, three strokes, the middle first.}

\textbf{The tall vertical} in the middle.

\textbf{A vertical down the left that turns} and runs right along the bottom.

\textbf{A vertical down the right}, standing on the base.

\subsection*{Your turn}
\begin{itemize}
\raggedright
  \item \emph{Write it:} \zh{山}, all three, with the model in front of you
  \item \emph{Cover:} write it again from memory, and count the pieces, not the strokes
\end{itemize}

\begin{tcolorbox}[breakable,colback=teal!4,colframe=teal!35!black,title={Before you move on}]
How many strokes in \zh{山}? (\textbf{Three}.) Which stroke comes first? (\textbf{The middle one}.) Which character you know has it on top? (\textbf{\zh{岁}}.)

Source: \href{https://raw.githubusercontent.com/chanind/hanzi-writer-data/68d10a4b21150cae5e1ebbd223eed289cf32d90c/data/\%E5\%B1\%B1.json}{Hanzi Writer Data \zh{山}, pinned stroke paths and medians}.
\end{tcolorbox}
\section[\texorpdfstring{\zh{星星} (xīngxing)}{星星 (xīngxing)}]{\zh{星星} (xīngxing) — star}

\label{lesson:ZH-C68-xingxing}



\begin{quote}
Before the new one: say the Mandarin for fish, then the Mandarin for mountain.
\end{quote}

\subsection*{You'll want to know: \zh{星星}}
\textbf{\zh{星星}} — \emph{xīngxing} — \textquotedblleft{}star\textquotedblright{}.

\begin{quote}
\textbf{\zh{天气很好晚上有星星}} \emph{tiānqì hěn hǎo wǎnshang yǒu xīngxing} — Fine weather: there'll be stars tonight.
\end{quote}

\begin{cousinweb}
\textbf{\zh{星星}} (\emph{xīngxing}) is built from \zh{星} 'star' said twice, the second time light. \zh{星} is inherited from Middle Chinese \emph{seng} (Guangyun, in Baxter's transcription).
\end{cousinweb}

\begin{grammarlens}[title={What you've built}]
One more word for everyday answers.
\end{grammarlens}

\subsection*{Your turn}
\begin{itemize}
\raggedright
  \item \emph{Say it:} \emph{xīngxing}
  \item \emph{Say it:} \emph{xīngxing}, once more
  \item \emph{Say it:} \emph{tiānqì hěn hǎo wǎnshang yǒu xīngxing} — Fine weather: there'll be stars tonight.
  \item \emph{Recall:} say the Mandarin for fish, then the Mandarin for mountain, then say \emph{xīngxing} again
\end{itemize}

\begin{tcolorbox}[breakable,colback=teal!4,colframe=teal!35!black,title={Before you move on}]
What does \zh{星星} mean? (\textquotedblleft{}Star\textquotedblright{}.) Say it once more.
\end{tcolorbox}
\section[cǎo]{cǎo — grass}

\label{lesson:ZH-C68-cao}



\begin{quote}
Before the new one: say the Mandarin for mountain, then the Mandarin for star.
\end{quote}

\subsection*{You'll want to know: cǎo}
\emph{cǎo} — \textquotedblleft{}grass\textquotedblright{}.

\begin{quote}
\emph{zhèr de cǎo hěn gāo} — The grass here is tall.
\end{quote}

Its characters come later; for now it is a word for the ear.

\begin{cousinweb}
\emph{cǎo} is inherited from Middle Chinese \emph{tshawX} (Guangyun, in Baxter's transcription), glossed, through an older book, as all the plants that grow.
\end{cousinweb}

\begin{grammarlens}[title={What you've built}]
One more word for everyday answers.
\end{grammarlens}

\subsection*{Your turn}
\begin{itemize}
\raggedright
  \item \emph{Say it:} \emph{cǎo}
  \item \emph{Say it:} \emph{cǎo}, once more
  \item \emph{Say it:} \emph{zhèr de cǎo hěn gāo} — The grass here is tall.
  \item \emph{Recall:} say the Mandarin for mountain, then the Mandarin for star, then say \emph{cǎo} again
\end{itemize}

\begin{tcolorbox}[breakable,colback=teal!4,colframe=teal!35!black,title={Before you move on}]
What does cǎo mean? (\textquotedblleft{}Grass\textquotedblright{}.) Say it once more.
\end{tcolorbox}
\section[huā]{huā — flower}

\label{lesson:ZH-C68-hua}



\begin{quote}
Before the new one: say the Mandarin for star, then the Mandarin for grass.
\end{quote}

\subsection*{You'll want to know: huā}
\emph{huā} — \textquotedblleft{}flower\textquotedblright{}.

\begin{quote}
\emph{māma xǐhuan huā} — Mum loves flowers.
\end{quote}

Its characters come later; for now it is a word for the ear.

\begin{cousinweb}
\emph{huā} is inherited from Middle Chinese \emph{xwae} (Guangyun, in Baxter's transcription), which the dictionary notes as the everyday form, in common use, of an older word for blossom.
\end{cousinweb}

\begin{grammarlens}[title={What you've built}]
One more word for everyday answers.
\end{grammarlens}

\subsection*{Your turn}
\begin{itemize}
\raggedright
  \item \emph{Say it:} \emph{huā}
  \item \emph{Say it:} \emph{huā}, once more
  \item \emph{Say it:} \emph{māma xǐhuan huā} — Mum loves flowers.
  \item \emph{Recall:} say the Mandarin for star, then the Mandarin for grass, then say \emph{huā} again
\end{itemize}

\begin{tcolorbox}[breakable,colback=teal!4,colframe=teal!35!black,title={Before you move on}]
What does huā mean? (\textquotedblleft{}Flower\textquotedblright{}.) Say it once more.
\end{tcolorbox}
\section[tàiyang]{tàiyang — sun}

\label{lesson:ZH-C68-taiyang}



\begin{quote}
Before the new one: say the Mandarin for grass, then the Mandarin for flower.
\end{quote}

\subsection*{You'll want to know: tàiyang}
\emph{tàiyang} — \textquotedblleft{}sun\textquotedblright{}.

\begin{quote}
\emph{jīntiān de tàiyang hěn dà} — The sun is strong today.
\end{quote}

Its characters come later; for now it is a word for the ear.

\begin{cousinweb}
\emph{tàiyang} is built from tài 'great' and yang, both inherited from Middle Chinese, \emph{thajH} and \emph{yang} (Guangyun, in Baxter's transcription). The dictionary glosses the second as the yang of yin and yang, 'high and bright': the great bright one.
\end{cousinweb}

\begin{grammarlens}[title={What you've built}]
One more word for everyday answers.
\end{grammarlens}

\subsection*{Your turn}
\begin{itemize}
\raggedright
  \item \emph{Say it:} \emph{tàiyang}
  \item \emph{Say it:} \emph{tàiyang}, once more
  \item \emph{Say it:} \emph{jīntiān de tàiyang hěn dà} — The sun is strong today.
  \item \emph{Recall:} say the Mandarin for grass, then the Mandarin for flower, then say \emph{tàiyang} again
\end{itemize}

\begin{tcolorbox}[breakable,colback=teal!4,colframe=teal!35!black,title={Before you move on}]
What does tàiyang mean? (\textquotedblleft{}Sun\textquotedblright{}.) Say it once more.
\end{tcolorbox}
